\newpage
\section{Mục tiêu}

\subsection{Mục tiêu chung}

Đề tài hướng tới việc nghiên cứu tường tận và ứng dụng phương pháp Học máy nhân quả (Causal Machine Learning), trọng tâm là mô hình đồ thị Mạng Bayes (Bayesian Network), nhằm giải quyết triệt để sự đánh đổi gay gắt giữa hiệu suất và khả năng giải thích trong các hệ thống đánh giá rủi ro tín dụng hiện nay. Trên cơ sở đó, mục tiêu cốt lõi là thiết kế và hiện thực hóa một khung dự báo rủi ro hoàn chỉnh, không chỉ đảm bảo năng lực phân loại chính xác trước đặc thù phức tạp và phân phối lệch của dữ liệu tài chính, mà còn khai thác tối đa khả năng giải thích nội tại (intrinsic interpretability) của cấu trúc mạng lưới nhân quả. Vượt lên trên việc ứng dụng thuật toán thuần túy, hệ thống này tập trung vào việc trích xuất minh bạch các luồng tác động gốc rễ thông qua kỹ thuật suy luận đa chiều, từ đó tạo tiền đề xây dựng một công cụ Business Intelligence (BI) trực quan, bám sát nghiệp vụ thực tế nhằm hỗ trợ an toàn và tin cậy cho các quyết định phê duyệt tài chính.

\subsection{Đóng góp của đề tài}

Đóng góp cốt lõi của đề tài không chỉ dừng lại ở việc ứng dụng thành công một thuật toán dự đoán rủi ro, mà quan trọng hơn là việc hiện thực hóa và tối ưu một khung phương pháp luận Học máy nhân quả (Causal Machine Learning) chuyên sâu cho dữ liệu tài chính. Cụ thể, nghiên cứu đã giải quyết triệt để sự đánh đổi gay gắt giữa hiệu suất phân loại và khả năng diễn giải bằng cách thiết lập một quy trình học Mạng Bayes (Bayesian Network) toàn diện. Thông qua việc chứng minh tính vượt trội của cấu trúc đồ thị xác suất trong việc bảo toàn và minh bạch hóa các luồng tác động gốc rễ so với các phương pháp XAI hậu kiểm truyền thống (như LIME, SHAP), đề tài cung cấp một giải pháp thực tiễn có khả năng trực tiếp nắm bắt logic nhân quả. Từ đó, hệ thống loại bỏ tối đa hiện tượng đứt gãy thông tin hay việc học sai các tương quan nhiễu bề mặt—yếu tố sống còn đối với một hệ thống thẩm định tín dụng cần tuân thủ các quy định pháp lý khắt khe về minh bạch thuật toán như GDPR.

Về mặt thực tiễn, kết quả nghiên cứu được chuyển hóa thành một công cụ Business Intelligence (BI Dashboard) trực quan, đóng vai trò như một trợ lý phân tích rủi ro hoạt động với độ tin cậy cao, giúp cá nhân hóa quá trình chẩn đoán nguyên nhân vỡ nợ cho từng hồ sơ và hỗ trợ ra quyết định an toàn cho đội ngũ chuyên viên tín dụng. Vượt ra khỏi phạm vi một khóa luận, nghiên cứu còn đóng góp một phương pháp luận tham chiếu (khung 5 bước: S\_Pro, L\_Pro, K\_Pro, BN\_Pro và In\_Pro) cho việc giải quyết nút thắt về chi phí tính toán khổng lồ (bài toán NP-hard) khi học cấu trúc mạng trên các không gian dữ liệu mất cân bằng và phức tạp. Điều này chứng minh rằng việc kết hợp tư duy đồ thị nhân quả với sức mạnh học máy là hướng đi tất yếu để xử lý các mạng lưới dữ liệu nghiệp vụ chằng chịt, tạo tiền đề mở rộng và nhân rộng mô hình sang các bài toán quản trị rủi ro cấp doanh nghiệp quy mô lớn. 